\documentclass[12pt]{article}
\usepackage[utf8]{inputenc}
\usepackage[T1]{fontenc}
\usepackage[margin=1in]{geometry}
\usepackage{setspace}
\usepackage{xcolor}
\usepackage[colorlinks=true, urlcolor=blue!60!black, linkcolor=blue!60!black]{hyperref}
\setlength{\parindent}{2em}
\setlength{\parskip}{1.25em}
\renewcommand{\baselinestretch}{1.0}
\begin{document}
\begin{center}
{\LARGE \textit{Taylor v. United States}, 579 U.S. 301 (2016)\par}
\vspace{0.8em}
{\large Case Brief\par}
\end{center}
\vspace{1em}

\section{Memory Jogger}
Applying \textit{Gonzales v. Raich} in a 7--1 decision, the Supreme Court held that the government satisfies the Hobbs Act's interstate commerce element merely by proving that a defendant robbed or attempted to rob a drug dealer of drugs or drug proceeds, without proving that the targeted drugs crossed state lines.

\section{Detailed Case Facts}
Beginning as early as 2009, an outlaw gang known as the ``Southwest Goonz'' carried out a series of home-invasion robberies targeting drug dealers around Roanoke, Virginia, because dealers typically keep substantial cash and illegal drugs at home and are reluctant to report thefts to police (\href{https://www.law.cornell.edu/supremecourt/text/14-6166}{Cornell LII}; \href{https://www.courtlistener.com/opinion/3214883/taylor-v-united-states/}{CourtListener}). Petitioner David Anthony Taylor participated in two such home invasions, leading to federal charges under the Hobbs Act, 18 U.S.C.\ \S~1951(a)---which criminalizes robbery or attempted robbery that ``in any way or degree obstructs, delays, or affects commerce''---and under 18 U.S.C.\ \S~924(c) for using a firearm in furtherance of a crime of violence (\href{https://www.law.cornell.edu/supremecourt/text/14-6166}{Cornell LII}).

The first robbery occurred in August 2009 at the home of Josh Whorley, whom the gang believed dealt ``exotic and high grade'' marijuana (\href{https://www.law.cornell.edu/supremecourt/text/14-6166}{Cornell LII}). Taylor and accomplices broke into Whorley's home, assaulted Whorley and his girlfriend, and demanded to know where the drugs and money were; finding neither stash, they fled with jewelry, \$40, two cell phones, and a single marijuana cigarette (\href{https://www.law.cornell.edu/supremecourt/text/14-6166}{Cornell LII}). The second robbery occurred in October 2009 at the home of William Lynch after an informant reported having previously robbed Lynch of 20 pounds of marijuana outside his house (\href{https://www.law.cornell.edu/supremecourt/text/14-6166}{Cornell LII}). Taylor and others broke in, held Lynch's wife and young children at gunpoint, assaulted his wife, and demanded drugs and money, ultimately leaving with only a cell phone (\href{https://www.law.cornell.edu/supremecourt/text/14-6166}{Cornell LII}).

\section{Procedural History}
\begin{enumerate}
\item \textbf{U.S. District Court for the Western District of Virginia:} On July 26, 2012, Taylor was indicted on two counts of Hobbs Act robbery (18 U.S.C.\ \S~1951(a)) and firearm counts (\href{https://www.law.cornell.edu/supct/cert/14-6166}{Cornell LII Preview}; \href{https://www.law.cornell.edu/supremecourt/text/14-6166}{Cornell LII}). His first trial ended in a hung jury after Taylor argued that robbing locally grown marijuana did not affect interstate commerce (\href{https://www.law.cornell.edu/supremecourt/text/14-6166}{Cornell LII}; \href{https://en.wikipedia.org/wiki/Taylor_v._United_States_(2016)}{Case Overview}). At retrial, at the Government's request, the district court barred Taylor from introducing evidence that the targeted dealers dealt solely in locally grown marijuana and twice denied his motions for judgment of acquittal, ruling as a matter of law that proof of attempting to rob drug dealers satisfied the commerce element (\href{https://www.law.cornell.edu/supremecourt/text/14-6166}{Cornell LII}). The jury convicted Taylor on both Hobbs Act counts and one \S~924(c) firearm count, and he was sentenced to 28 years' imprisonment (\href{https://www.law.cornell.edu/supremecourt/text/14-6166}{Cornell LII}; \href{https://www.law.georgetown.edu/american-criminal-law-review/wp-content/uploads/sites/15/2023/02/53-0-Evola-you-shall-go-no-further.pdf}{American Criminal Law Review}).
\item \textbf{U.S. Court of Appeals for the Fourth Circuit:} On June 6, 2014, the Fourth Circuit affirmed (\textit{United States v. Taylor}, 754 F.3d 217 (4th Cir. 2014)), holding that because drug dealing in the aggregate necessarily affects interstate commerce, the Government needed only to prove that Taylor depleted or attempted to deplete the assets of a drug-dealing operation (\href{https://caselaw.findlaw.com/court/us-4th-circuit/1669252.html}{FindLaw Fourth Circuit}; \href{https://www.law.cornell.edu/supremecourt/text/14-6166}{Cornell LII}).
\item \textbf{Supreme Court of the United States:} The Court granted certiorari on October 1, 2015, to resolve a circuit split with the Second and Seventh Circuits (\textit{United States v. Needham}, 604 F.3d 673 (2d Cir. 2010); \textit{United States v. Peterson}, 236 F.3d 848 (7th Cir. 2001)) (\href{https://www.oyez.org/cases/2015/14-6166}{Oyez}; \href{https://www.law.cornell.edu/supremecourt/text/14-6166}{Cornell LII}). Argued February 23, 2016---before an eight-Justice Court following Justice Scalia's death---the case was decided June 20, 2016, affirming the Fourth Circuit 7--1 (\href{https://www.oyez.org/cases/2015/14-6166}{Oyez}; \href{https://www.courtlistener.com/opinion/3214883/taylor-v-united-states/}{CourtListener}).
\end{enumerate}

\section{Judicial Votes}
\begin{itemize}
\item \textbf{Vote Breakdown:} 7--1 affirming the Fourth Circuit (eight-Justice Court following the vacancy left by Justice Scalia's death) (\href{https://www.law.cornell.edu/supremecourt/text/14-6166}{Cornell LII}; \href{https://www.oyez.org/cases/2015/14-6166}{Oyez}).
\item \textbf{Majority Opinion (7):} Justice Alito delivered the opinion of the Court, joined by Chief Justice Roberts and Justices Kennedy, Ginsburg, Breyer, Sotomayor, and Kagan (\href{https://www.law.cornell.edu/supremecourt/text/14-6166}{Cornell LII}).
\item \textbf{Concurring Opinions:} None (\href{https://www.law.cornell.edu/supremecourt/text/14-6166}{Cornell LII}).
\item \textbf{Dissenting Opinion (1):} Justice Thomas filed a lone dissenting opinion (\href{https://www.law.cornell.edu/supremecourt/text/14-6166}{Cornell LII}).
\end{itemize}

\section{Holding}
\begin{itemize}
\item \textbf{Judgment:} Affirmed. The evidence that Taylor knowingly targeted drug dealers to steal drugs and drug proceeds was sufficient to satisfy the Hobbs Act's commerce element (\href{https://www.law.cornell.edu/supremecourt/text/14-6166}{Cornell LII}).
\item \textbf{Rule Announced:} In a Hobbs Act prosecution for robbery or attempted robbery targeting a drug dealer, the Government need not prove that the drugs traveled or were destined to travel across state lines; proving beyond a reasonable doubt that the defendant knowingly stole or attempted to steal drugs or drug proceeds from a drug dealer satisfies the Act's commerce element as a matter of law, because under \textit{Gonzales v. Raich}, 545 U.S. 1 (2005), the entire market for illegal drugs---including purely intrastate production and sale---constitutes ``commerce over which the United States has jurisdiction'' under 18 U.S.C.\ \S~1951(b)(3) (\href{https://www.law.cornell.edu/supremecourt/text/14-6166}{Cornell LII}).
\item \textbf{Holding vs.\ Dicta:} The holding is expressly confined to cases in which a defendant targets drug dealers to steal drugs or drug proceeds (\href{https://www.law.cornell.edu/supremecourt/text/14-6166}{Cornell LII}; \href{https://www.law.cornell.edu/constitution-conan/article-1/section-8/clause-3/limits-on-federal-regulation-of-intrastate-activity}{Constitution Annotated}). The Court explicitly left unresolved what the Government must prove when a Hobbs Act robbery targets another type of business or victim, and its remarks that \textit{Taylor} merely applies without expanding \textit{Raich}---noting that the Court was not asked to reconsider \textit{Raich}---constitute limiting commentary and dicta rather than a reaffirmation of \textit{Raich} against direct constitutional challenge (\href{https://www.law.cornell.edu/supremecourt/text/14-6166}{Cornell LII}; \href{https://www.nationalreview.com/bench-memos/loose-evidentiary-ends-raich-taylor-v-united-states}{National Review}).
\end{itemize}

\section{Analysis of Opinions}

\subsection{Majority Opinion of Justice Alito (Joined by Chief Justice Roberts and Justices Kennedy, Ginsburg, Breyer, Sotomayor, and Kagan)}
Justice Alito framed the case as a ``straightforward'' application of existing statutory and Commerce Clause precedent (\href{https://www.law.cornell.edu/supremecourt/text/14-6166}{Cornell LII}). First, examining the text of the Hobbs Act, 18 U.S.C.\ \S\S~1951(a) and 1951(b)(3), the Court emphasized that the statute's language is ``unmistakably broad'' and reaches any obstruction, delay, or effect on ``all other commerce over which the United States has jurisdiction''---exercising the full constitutional reach of Congress's commerce power (\textit{Stirone v. United States}, 361 U.S. 212, 215 (1960); \textit{United States v. Culbert}, 435 U.S. 371, 373 (1978)) (\href{https://www.law.cornell.edu/supremecourt/text/14-6166}{Cornell LII}).

Second, turning to Commerce Clause doctrine under \textit{United States v. Lopez}, 514 U.S. 549 (1995), \textit{Wickard v. Filburn}, 317 U.S. 111 (1942), and \textit{United States v. Morrison}, 529 U.S. 598 (2000), the Court reiterated that Congress may regulate purely local economic activities that belong to an economic ``class of activities'' substantially affecting interstate commerce in the aggregate (\href{https://www.law.cornell.edu/supremecourt/text/14-6166}{Cornell LII}). Because \textit{Gonzales v. Raich}, 545 U.S. 1 (2005), held that the production, possession, and distribution of controlled substances---including purely intrastate marijuana---constitutes such a class of economic activities, the market for marijuana is ``commerce over which the United States has jurisdiction'' (\href{https://www.law.cornell.edu/supremecourt/text/14-6166}{Cornell LII}). Simply grafting \textit{Raich} onto the Hobbs Act, a robber who targets a marijuana dealer's drugs or proceeds necessarily affects or attempts to affect federally regulable commerce (\href{https://www.law.cornell.edu/supremecourt/text/14-6166}{Cornell LII}).

Third, the Court rejected Taylor's argument that because the Hobbs Act contains an express jurisdictional element (unlike the Controlled Substances Act in \textit{Raich}), the Government must prove beyond a reasonable doubt that the specific drugs crossed state lines or that the victim ran an interstate enterprise (\href{https://www.law.cornell.edu/supremecourt/text/14-6166}{Cornell LII}). Justice Alito explained that Taylor confused the \textit{standard of proof} with the \textit{legal meaning of the element}: while the jury must find the facts beyond a reasonable doubt, the meaning of ``commerce over which the United States has jurisdiction'' is a question of law already settled by \textit{Raich}, and under \textit{Perez v. United States}, 402 U.S. 146 (1971), a minimal or trivial individual effect within a regulated class suffices (\href{https://www.law.cornell.edu/supremecourt/text/14-6166}{Cornell LII}). Finally, the Court emphasized that Taylor had not asked the Court to overrule \textit{Raich} and explicitly limited its holding to robberies targeting drug dealers' drugs or drug proceeds (\href{https://www.law.cornell.edu/supremecourt/text/14-6166}{Cornell LII}).

\subsection{Concurring Opinions}
No separate concurring opinions were filed (\href{https://www.law.cornell.edu/supremecourt/text/14-6166}{Cornell LII}).

\subsection{Dissenting Opinion of Justice Thomas}
Justice Thomas dissented alone, advancing three interconnected constitutional, doctrinal, and criminal-procedure objections (\href{https://www.law.cornell.edu/supremecourt/text/14-6166}{Cornell LII}).

First, grounded in the original meaning of the Constitution, Justice Thomas emphasized that Congress lacks a general police power over local crime; the Constitution delegates federal criminal authority over only four enumerated offenses (counterfeiting, piracies and felonies on the high seas, offenses against the law of nations, and treason) (\href{https://www.law.cornell.edu/supremecourt/text/14-6166}{Cornell LII}). Because robbery is not ``Commerce''---which at the Founding meant buying, selling, bartering, and transporting goods across state lines---Congress may punish robbery only when ``necessary and proper'' to carrying out its interstate commerce power, requiring an ``obvious, simple, and direct relation'' to interstate commerce (\href{https://www.law.cornell.edu/supremecourt/text/14-6166}{Cornell LII}). Under that constitutional limit, the Hobbs Act validly reaches robberies affecting channels or instrumentalities of interstate commerce (such as shutting down an interstate highway or robbing an interstate truck driver), but not purely local home invasions (\href{https://www.law.cornell.edu/supremecourt/text/14-6166}{Cornell LII}).

Second, even under modern precedent, Justice Thomas argued that applying the ``substantial-effects'' aggregation test to the Hobbs Act has ``no stopping point'' and obliterates the distinction between ``what is truly national and what is truly local'' (\href{https://www.law.cornell.edu/supremecourt/text/14-6166}{Cornell LII}). Moreover, \textit{Taylor} actually \textit{expanded} \textit{Raich}: \textit{Raich} held only that Congress had a deferential ``rational basis'' to believe intrastate marijuana affected the interstate market, not that every robbery targeting a dealer in fact affects commerce beyond a reasonable doubt (\href{https://www.law.cornell.edu/supremecourt/text/14-6166}{Cornell LII}; \href{https://www.nationalreview.com/bench-memos/loose-evidentiary-ends-raich-taylor-v-united-states}{National Review}).

Third, Justice Thomas warned that the majority eroded foundational criminal-law protections---including the Sixth Amendment and Due Process requirement of proof beyond a reasonable doubt of every element (\textit{In re Winship}, 397 U.S. 358 (1970)), the rule of lenity (\textit{United States v. Wiltberger}, 5 Wheat.\ 76 (1820)), and the clear-statement federalism canon (\textit{Jones v. United States}, 529 U.S. 848 (2000))---by relieving prosecutors of proving any actual nexus to interstate commerce (\href{https://www.law.cornell.edu/supremecourt/text/14-6166}{Cornell LII}).

\section{Examples: Future Applications}
\begin{enumerate}
\item \textbf{Same-Side (Hobbs Act Commerce Element Satisfied):} A defendant breaks into a local methamphetamine dealer's apartment in rural Ohio intending to steal cash proceeds from intrastate sales, even though the dealer manufactures the methamphetamine locally using locally purchased supplies and sells only to neighbors. Under \textit{Taylor}, because the defendant targeted a controlled-substance dealer's drugs or proceeds, the commerce element is satisfied as a matter of law (\href{https://www.law.cornell.edu/supremecourt/text/14-6166}{Cornell LII}).
\item \textbf{Same-Side (Hobbs Act Commerce Element Satisfied):} Two assailants rob an unlicensed local street-level cocaine seller at gunpoint, demanding his stash and cash, but the seller happens to have already sold out for the night and carries only \$15. Just as in \textit{Taylor}, an attempted robbery targeting a drug dealer's drugs or proceeds satisfies the Hobbs Act even when no drugs are recovered (\href{https://www.law.cornell.edu/supremecourt/text/14-6166}{Cornell LII}).
\item \textbf{Opposite-Side (Hobbs Act Commerce Element Not Automatically Satisfied):} A burglar breaks into a retired schoolteacher's private home and steals a personal heirloom watch and \$200 in household cash, with no connection to drug trafficking or any commercial enterprise. Because \textit{Taylor} is expressly limited to robberies targeting drug dealers' drugs or proceeds, the Government cannot rely on \textit{Taylor}'s categorical rule and must prove an actual nexus to interstate commerce (\href{https://www.law.cornell.edu/supremecourt/text/14-6166}{Cornell LII}; \href{https://www.law.cornell.edu/wex/Hobbs_Act}{Cornell Wex}).
\item \textbf{Opposite-Side (Hobbs Act Commerce Element Not Automatically Satisfied):} A robber mugging a pedestrian on a residential street steals a personal bicycle and a wallet from someone who happens to smoke marijuana recreationally, without targeting any drug-dealing enterprise or intending to steal drugs or drug proceeds. Because the victim was not targeted as a drug dealer for drugs or drug proceeds, \textit{Taylor}'s categorical rule does not apply (\href{https://www.law.cornell.edu/supremecourt/text/14-6166}{Cornell LII}).
\item \textbf{Fence-Sitter (Close Call):} A robber holds up a state-licensed, strictly intrastate recreational cannabis dispensary in Colorado (or robs a homeowner who grows small-scale wheat solely for personal consumption, a hypothetical raised by Justice Kennedy at oral argument). While \textit{Taylor} treats illegal drug markets as categorically covered under \textit{Raich}, courts face contested questions over whether \textit{Taylor}'s categorical shortcut extends to state-legalized intrastate cannabis businesses or to non-drug regulated commodities like homegrown agricultural goods (\href{https://fedsoc.org/commentary/fedsoc-blog/taylor-v-united-states}{Federalist Society}; \href{https://www.scotusblog.com/2025/10/the-case-that-made-federalism-go-up-in-smoke/}{SCOTUSblog}).
\end{enumerate}

\section{Critique}
\textbf{Originalist and Textualist Perspectives:} Originalist and federalism-oriented commentators strongly echo Justice Thomas's dissent, criticizing \textit{Taylor} for entrenching \textit{Raich}'s sweeping aggregation doctrine in federal criminal law and effectively converting a local home invasion into a 20-year federal felony (\href{https://www.nationalreview.com/bench-memos/loose-evidentiary-ends-raich-taylor-v-united-states}{National Review}; \href{https://www.scotusblog.com/2025/10/the-case-that-made-federalism-go-up-in-smoke/}{SCOTUSblog}). Conservative analysts note that Justice Alito's opinion hinted at discomfort with \textit{Raich}---stressing that the Court had not been asked to reconsider it---while nonetheless allowing \textit{Raich}'s deferential rational-basis test to displace proof beyond a reasonable doubt of a statutory jurisdictional element (\href{https://www.nationalreview.com/bench-memos/loose-evidentiary-ends-raich-taylor-v-united-states}{National Review}). Commentators also point out that the Hobbs Act was enacted in 1946 to combat labor-related ``highway robbery'' and extortion obstructing interstate shipments, not purely local residential thefts (\href{https://www.law.georgetown.edu/american-criminal-law-review/wp-content/uploads/sites/15/2023/02/53-0-Evola-you-shall-go-no-further.pdf}{American Criminal Law Review}).

\textbf{Progressive and Doctrinal Perspectives:} While progressive constitutional theorists generally defend broad post-New Deal Commerce Clause authority to address national economic markets, progressive criminal-justice reformers and defense advocates sharply criticized \textit{Taylor} for accelerating federal overcriminalization and mass incarceration (\href{https://thehill.com/homenews/284991-supreme-courts-hobbs-act-ruling-means-increased-mass-incarceration/}{The Hill}; \href{https://www.law.georgetown.edu/american-criminal-law-review/wp-content/uploads/sites/15/2023/02/53-0-Evola-you-shall-go-no-further.pdf}{American Criminal Law Review}). Critics highlight that Taylor---described by government counsel at oral argument as a ``bit of a tagalong'' who stole \$40 and cell phones---received a 28-year federal prison sentence for conduct already punishable under Virginia state law (\href{https://www.law.georgetown.edu/american-criminal-law-review/wp-content/uploads/sites/15/2023/02/53-0-Evola-you-shall-go-no-further.pdf}{American Criminal Law Review}; \href{https://thehill.com/homenews/284991-supreme-courts-hobbs-act-ruling-means-increased-mass-incarceration/}{The Hill}). Doctrinal scholars further question the logical tension between \textit{Taylor}'s insistence that it does not expand \textit{Raich} and its use of a legislative-deference precedent to conclusively establish an element of a federal crime (\href{https://www.nationalreview.com/bench-memos/loose-evidentiary-ends-raich-taylor-v-united-states}{National Review}).

\textbf{Underlying Competing Values:} \textit{Taylor} pits doctrinal coherence and effective federal law enforcement against enumerated-powers federalism and criminal-defendant safeguards. On one side are stare decisis (taking \textit{Raich} and \textit{Stirone} at face value) and the federal government's interest in deploying unified national tools against violent gangs preying on illicit drug markets. On the other side are constitutional limits preventing a general federal police power, respect for state sovereignty over local violent crime, and the Sixth Amendment guarantee that prosecutors must prove every statutory element to a jury beyond a reasonable doubt (\href{https://www.law.cornell.edu/supremecourt/text/14-6166}{Cornell LII}; \href{https://www.law.georgetown.edu/american-criminal-law-review/wp-content/uploads/sites/15/2023/02/53-0-Evola-you-shall-go-no-further.pdf}{American Criminal Law Review}).

\section{Key Quotations}
\begin{quote}
``The case now before us requires no more than that we graft our holding in \textit{Raich} onto the commerce element of the Hobbs Act.''
\end{quote}
\noindent --- Justice Alito, opinion of the Court, 579 U.S. at 307 (\href{https://www.law.cornell.edu/supremecourt/text/14-6166}{Cornell LII}; \href{https://www.courtlistener.com/opinion/3214883/taylor-v-united-states/}{CourtListener}).

\begin{quote}
``Rather, to satisfy the Act's commerce element, it is enough that a defendant knowingly stole or attempted to steal drugs or drug proceeds, for, as a matter of law, the market for illegal drugs is `commerce over which the United States has jurisdiction.' ''
\end{quote}
\noindent --- Justice Alito, opinion of the Court, 579 U.S. at 309 (\href{https://www.law.cornell.edu/supremecourt/text/14-6166}{Cornell LII}).

\begin{quote}
``Our holding today is limited to cases in which the defendant targets drug dealers for the purpose of stealing drugs or drug proceeds.''
\end{quote}
\noindent --- Justice Alito, opinion of the Court, 579 U.S. at 310 (\href{https://www.law.cornell.edu/supremecourt/text/14-6166}{Cornell LII}).

\begin{quote}
``Under the Court's decision today, the Government can obtain a Hobbs Act conviction without proving that the defendant's robbery in fact affected interstate commerce---or any commerce.''
\end{quote}
\noindent --- Justice Thomas, dissenting, 579 U.S. at 311 (\href{https://www.law.cornell.edu/supremecourt/text/14-6166}{Cornell LII}).

\begin{quote}
``By applying the substantial-effects test to the criminal prohibition before us, the Court effectively gives Congress a police power.''
\end{quote}
\noindent --- Justice Thomas, dissenting, 579 U.S. at 319 (\href{https://www.law.cornell.edu/supremecourt/text/14-6166}{Cornell LII}).

\end{document}
